\documentclass[12pt]{article}
\usepackage[T1]{fontenc}
\usepackage[utf8]{inputenc}
\usepackage{lmodern}
\usepackage{textcomp}
\usepackage{enumitem}
\usepackage{geometry}
\geometry{margin=1in}
\begin{document}
\begin{center}
Answer Key - Economics - K-12 Level Assessment 2 \
\small (Answers are ordered exactly as in the question paper.)
\end{center}

\section*{Multiple Choice}
\begin{enumerate}[label=\arabic*.]
\item \textbf{Answer:} D
\\ \textit{Options:} A) have not changed significantly.; B) are perfectly constant.; C) have decreased significantly.; D) have increased significantly.
\\ \textit{Note:} The correct answer is supported by the observation that both the money supply (V) and the quantity of goods and services (Q) in the economy have experienced significant growth due to factors such as technological advancements, population growth, and increased consumer demand over time.
\item \textbf{Answer:} D
\\ \textit{Options:} A) Both the supply of and the demand for the good decrease.; B) Both the supply of and the demand for the good increase.; C) The supply increases and the demand for the good decreases.; D) None of the above.
\\ \textit{Note:} The correct answer is "None of the above" because in economics, price and quantity changes for a particular good cannot be determined simultaneously without additional information about the market conditions and demand and supply shifts.
\item \textbf{Answer:} A
\\ \textit{Options:} A) A decrease in the required reserve ratio; B) An increase in the discount rate; C) The selling of bonds by the Federal Reserve; D) An increase in the fraction of deposits that must be held by banks
\\ \textit{Note:} A decrease in the required reserve ratio allows banks to hold less money in reserve and lend more, which increases the overall money supply in the economy.
\item \textbf{Answer:} D
\\ \textit{Options:} A) The price level rises but real GDP falls.; B) Both the price level and real GDP rise.; C) The price level rises but the change in real GDP is uncertain.; D) Real GDP rises but the change in the price level is uncertain.
\\ \textit{Note:} When both short-run aggregate supply and aggregate demand increase, real GDP will definitely rise due to the higher production and consumption levels, but the change in the price level is uncertain because the extent of the shifts in supply and demand can lead to varying impacts on prices.
\item \textbf{Answer:} C
\\ \textit{Options:} A) a persistent shortage of the good.; B) an increase in total welfare.; C) a persistent surplus of the good.; D) elimination of deadweight loss.
\\ \textit{Note:} A price floor set above the market equilibrium price leads to a situation where the quantity supplied exceeds the quantity demanded, resulting in a persistent surplus of the good.
\end{enumerate}

\section*{True / False}
\begin{enumerate}[label=\arabic*.]
\setcounter{enumi}{5}
\item \textbf{Answer:} True
\\ \textit{Options:} A) True; B) False
\\ \textit{Note:} The statement is true because a demand price elasticity of 0.78 indicates inelastic demand, meaning that consumers will not significantly reduce their quantity demanded in response to a price increase, allowing the business to raise prices and increase total revenue.
\item \textbf{Answer:} False
\\ \textit{Options:} A) True; B) False
\\ \textit{Note:} The statement is false because average total cost is calculated by dividing total costs (which include both fixed and variable costs) by the number of units produced, not just by total fixed costs.
\item \textbf{Answer:} True
\\ \textit{Options:} A) True; B) False
\\ \textit{Note:} Automatic stabilizers, such as unemployment benefits and progressive tax systems, automatically increase government spending or reduce taxes during a recession, which helps to cushion the economic downturn and support consumer spending.
\item \textbf{Answer:} True
\\ \textit{Options:} A) True; B) False
\\ \textit{Note:} Oligopoly is defined as a market structure where a small number of firms hold a significant market share, leading to limited competition and interdependent decision-making among the dominant producers.
\item \textbf{Answer:} False
\\ \textit{Options:} A) True; B) False
\\ \textit{Note:} The statement is false because an 8 percent increase in nominal income, when adjusted for a 2 percent rise in the CPI, results in a real income increase of approximately 5.88 percent, not a decrease.
\end{enumerate}

\section*{Long Form Response}
\begin{enumerate}[label=\arabic*.]
\setcounter{enumi}{10}
\item \textbf{Answer:} Hiring additional labor for a firm that pays $15 per worker and sells products for $3 each has important implications, especially when considering the marginal product of the third worker. If the third worker's contribution leads to the production of more goods, the firm can sell these additional products, increasing overall revenue. However, as more workers are hired, the marginal revenue product of labor (MRPL) tends to decrease because each new worker may contribute less to total output than the previous one. Therefore, the firm should ensure that it continues to hire workers as long as the additional revenue generated by each new worker exceeds their wage cost. This careful balance helps the firm maximize its profits while managing costs effectively.
\\ \textit{Note:} correct because it illustrates how hiring additional labor can initially boost revenue through increased production, but also highlights the potential diminishing returns on each worker's contribution to output, affecting the firm's overall profitability.
\item \textbf{Answer:} An increase in the price of tubas can lead to higher wages for tuba makers due to changes in supply and demand. When the price of tubas rises, it often indicates a higher demand for these instruments, which can encourage tuba manufacturers to produce more. To meet this increased production, companies may need to hire more workers or pay existing workers higher wages to keep up with the demand. Additionally, if the cost of materials and labor remains steady while the prices of tubas increase, the profit margins for tuba makers could expand, allowing them to offer better compensation to their employees. Overall, the dynamics of the market can create a positive feedback loop, where increased prices result in higher wages for those involved in making tubas.
\\ \textit{Note:} correct because an increase in the price of tubas typically signals higher demand, prompting manufacturers to potentially raise wages to attract more workers or incentivize current employees to meet the increased production needs.
\end{enumerate}

\vfill
\noindent\textit{End of Answer Key}
\end{document}